%!TEX program = pdflatex
%
% ECG SIMULATOR — COMPREHENSIVE TECHNICAL DOCUMENTATION
% 
% Complete, rigorous technical reference for simultaneous multi-audience consumption:
% • Physicians and medical students (no programming/math background assumed)
% • Bioscientists and physiologists (limited programming/physics background)
% • Mathematicians, physicists, and software engineers (full rigor required)
%
% All equations in correct LaTeX math notation.
% All Python code with exact parameters and dependencies.
% Exhaustive GUI reference: every parameter, every range, every meaning.
% Complete pathology catalog with step-by-step GUI reproduction.
% Architecture diagrams in TikZ.
% Publication-standard formatting: zero layout defects.
%
% Compilation: pdflatex -interaction=nonstopmode ECG_Simulator_Comprehensive.tex (twice for TOC)
%
% ============================================================================
% PREAMBLE — REQUIRED PACKAGES
% ============================================================================

\documentclass[12pt,oneside,a4paper]{book}

% ── Geometry ────────────────────────────────────────────────────────────
% Page layout: 1 in margins all around, generous spacing
\usepackage[
  top=1in,
  bottom=1in,
  left=1in,
  right=1in,
  headheight=15pt,
  headsep=0.25in,
  footskip=0.5in
]{geometry}

% ── Fonts & Encoding ────────────────────────────────────────────────────
\usepackage[utf-8]{inputenc}
\usepackage[T1]{fontenc}
\usepackage{lmodern}                    % Latin Modern (scalable, professional)

% ── Math & Symbols ─────────────────────────────────────────────────────
\usepackage{amsmath}                    % Core math environments
\usepackage{amssymb}                    % Mathematical symbols (ℝ, ℂ, ∀, ∃, ∈, etc.)
\usepackage{mathtools}                  % Extended math tools (equations, aligned environments)
\usepackage{array}                      % Enhanced array/table environments

% ── Colors ──────────────────────────────────────────────────────────────
\usepackage{xcolor}                     % Named colors, custom definitions

% ── Graphics & TikZ ─────────────────────────────────────────────────────
\usepackage{graphicx}                   % Include graphics
\usepackage{tikz}                       % Vector diagrams (class hierarchy, data flow, etc.)
\usepackage{pgfplots}                   % Numerical plotting on LaTeX canvas
\pgfplotsset{compat=1.18}               % Use pgfplots compatibility with recent versions
\usetikzlibrary{shapes,arrows,positioning,calc,backgrounds,decorations.pathreplacing}

% ── Floats & Captions ───────────────────────────────────────────────────
\usepackage{float}                      % Better float control ([H] placement)
\usepackage{subcaption}                 % Subfigures and subtables
\usepackage{caption}                    % Enhanced caption formatting
\captionsetup{
  font=small,
  labelfont=bf,
  justification=raggedright,
  singlelinecheck=false
}

% ── Tables ──────────────────────────────────────────────────────────────
\usepackage{booktabs}                   % Professional table formatting (\toprule, \midrule, \bottomrule)
\usepackage{tabularx}                   % Table with fixed width and auto-wrapping columns
\usepackage{multirow}                   % Cells spanning multiple rows
\usepackage{array}                      % Enhanced column specifications

% ── Code Listings ──────────────────────────────────────────────────────
\usepackage{listings}                   % Code syntax highlighting
\lstset{
  language=Python,
  basicstyle=\ttfamily\small,
  keywordstyle=\bfseries\color{blue},
  commentstyle=\itshape\color{gray},
  stringstyle=\color{red},
  showstringspaces=false,
  breaklines=true,
  breakatwhitespace=true,
  numbers=left,
  numberstyle=\tiny\color{gray},
  numbersep=5pt,
  frame=single,
  frameround=tttt,
  rulecolor=\color{lightgray},
  backgroundcolor=\color{white},
  framerule=0.5pt,
  xleftmargin=0.5cm,
  xrightmargin=0.5cm,
  aboveskip=0.5cm,
  belowskip=0.5cm,
}

% ── Hyperlinks & References ────────────────────────────────────────────
\usepackage{hyperref}                   % Cross-references, TOC links, URLs
\hypersetup{
  colorlinks=true,
  linkcolor=blue,
  urlcolor=blue,
  citecolor=black,
  pdfauthor={ECG Simulator Project},
  pdftitle={ECG Simulator -- Comprehensive Technical Documentation},
  pdfsubject={Cardiac Electrophysiology Simulation},
  bookmarksnumbered=true,
  bookmarksopen=false,
}

% ── Headers & Footers ──────────────────────────────────────────────────
\usepackage{fancyhdr}
\pagestyle{fancy}
\fancyhf{}
\fancyhead[L]{\nouppercase{\leftmark}}
\fancyhead[R]{\thepage}
\renewcommand{\headrulewidth}{0.4pt}
\renewcommand{\footrulewidth}{0pt}

% ── Sections & Margins ─────────────────────────────────────────────────
\usepackage{titlesec}                   % Customize section formatting
\titleformat{\chapter}[display]
  {\normalfont\huge\bfseries}
  {\chaptertitlename\ \thechapter}{20pt}{\huge}
\titleformat{\section}
  {\normalfont\Large\bfseries}{\thesection}{1em}{}
\titleformat{\subsection}
  {\normalfont\large\bfseries}{\thesubsection}{1em}{}
\titleformat{\subsubsection}
  {\normalfont\normalsize\bfseries}{\thesubsubsection}{1em}{}

% ── Paragraphing ───────────────────────────────────────────────────────
\setlength{\parskip}{0.5em}
\setlength{\parindent}{0.5in}

% ── Lists ──────────────────────────────────────────────────────────────
\usepackage{enumitem}                   % Customizable lists
\setlist[enumerate]{leftmargin=1cm, label=\arabic*.}
\setlist[itemize]{leftmargin=0.5cm, label=\textbullet}

% ── Miscellaneous ──────────────────────────────────────────────────────
\usepackage{fancyvrb}                   % Better verbatim environments
\usepackage{afterpage}                  % Cleaner page manipulation
\usepackage[hidelinks]{hyperref}        % Hyperlinks without visible boxes (reloaded for safety)

% ============================================================================
% CUSTOM COMMANDS & MACROS
% ============================================================================

% Inline code
\newcommand{\code}[1]{\texttt{#1}}

% Inline math (with proper spacing)
\newcommand{\math}[1]{$#1$}

% Parameter specification (name, unit, range)
\newcommand{\param}[4]{\textbf{\code{#1}} \\ \quad Type: \code{#2} \quad Unit: \emph{#3} \quad Range: $#4$}

% File/class reference
\newcommand{\file}[1]{\texttt{#1}}
\newcommand{\class}[1]{\textbf{\code{#1}}}

% Important notes
\newcommand{\note}[1]{\textbf{Note:} #1}
\newcommand{\warning}[1]{\textbf{\color{red}Warning:} #1}

% Key concept (first mention of important term)
\newcommand{\concept}[1]{\textbf{\emph{#1}}}

% ============================================================================
% FRONT MATTER
% ============================================================================

\title{%
  \Huge\textbf{ECG Simulator}\\
  \huge\textit{Comprehensive Technical Documentation}\\
  \large{A Multi-Phase Cardiac Electrophysiology Simulation Platform}
}

\author{%
  Project Documentation\\
  \textit{Complete technical reference for physicians, bioscientists, engineers}
}

\date{\today}

% ============================================================================
% BEGIN DOCUMENT
% ============================================================================

\begin{document}

% ── Title Page ──────────────────────────────────────────────────────────
\maketitle

% ── Table of Contents ───────────────────────────────────────────────────
\tableofcontents
\newpage

% ── List of Figures ─────────────────────────────────────────────────────
\listoffigures
\newpage

% ── List of Tables ──────────────────────────────────────────────────────
\listoftables
\newpage

% ============================================================================
% CHAPTER 1: PROJECT OVERVIEW & ARCHITECTURE
% ============================================================================

\chapter{Project Overview \& Architecture}

\section{1.1 Executive Summary}

The ECG Simulator is a Python-based cardiac electrophysiology simulation platform that generates clinically realistic 12-lead electrocardiograms (ECGs) under physiological and pathological conditions. The simulator spans three phases of increasing biological and computational sophistication:

\begin{itemize}
  \item \textbf{Phase 1/2}: Discrete conduction graph with hardcoded timing, suitable for fast prototyping and all conduction-level pathologies.
  \item \textbf{Phase 3}: One-dimensional monodomain cable with FitzHugh-Nagumo cellular electrophysiology, producing emergent refractory periods and wave-based QRS morphology.
  \item \textbf{Phase 4+}: Full 3D tissue modeling, advanced ion-channel models, and specialized arrhythmia mechanisms (planned).
\end{itemize}

The simulator is designed for \emph{simultaneous accessibility} to four distinct user groups:

\begin{enumerate}
  \item \textbf{Physicians \& Medical Students}: Understand normal and abnormal ECG patterns without programming knowledge.
  \item \textbf{Bioscientists \& Physiologists}: Study cardiac electrophysiology in a controlled computational environment.
  \item \textbf{Mathematicians \& Physicists}: Access rigorous mathematical models, numerical methods, and validation results.
  \item \textbf{Software Engineers \& Developers}: Extend and maintain the codebase, add new pathologies, and integrate with other systems.
\end{enumerate}

\section{1.2 Key Features}

\subsection{1.2.1 Clinical Accuracy}

\begin{itemize}
  \item 12-lead ECG standard output with clinically validated lead vectors.
  \item Normal sinus rhythm parameters calibrated to clinical targets:
    \begin{itemize}
      \item P-wave duration: 64 ms (clinical: 40--100 ms)
      \item PR interval: 196 ms (clinical: 120--200 ms)
      \item QRS duration: 70 ms (clinical: \textless 120 ms)
      \item Heart rate: adjustable 30--220 bpm
    \end{itemize}
  \item 15+ pathological patterns with exact clinical morphologies.
\end{itemize}

\subsection{1.2.2 Computational Performance}

\begin{itemize}
  \item Phase 1/2: \textgreater 1000 Hz output rate (highly efficient for real-time visualization).
  \item Phase 3: 2.2$\times$ real-time on contemporary hardware using NumPy fallback (Numba JIT optimization available when NumPy $\leq$ 2.4 is used).
  \item Continuous 24-hour simulation with bounded memory footprint (0.85 s beat retention window).
\end{itemize}

\subsection{1.2.3 Extensibility}

\begin{itemize}
  \item Plugin architecture for static and stateful pathologies.
  \item Runtime engine selection (Phase 1/2 vs Phase 3).
  \item Parameterized cell models (FitzHugh-Nagumo, Aliev-Panfilov).
  \item Single-equivalent-dipole forward model (Phase 5 to include boundary-element torso model).
\end{itemize}

\section{1.3 System Architecture}

\subsection{1.3.1 Layered Design}

The codebase follows a strict three-layer architecture:

\begin{table}[H]
\centering
\begin{tabularx}{\textwidth}{|l|X|X|}
\hline
\textbf{Layer} & \textbf{Responsibility} & \textbf{Key Modules} \\
\hline
\textbf{GUI / UI} & User interaction, parameter editing, ECG display, pathology selection &
  \file{main\_window.py}, \file{parameter\_panel.py}, \file{pathology\_panel.py}, \file{ecg\_display.py} \\
\hline
\textbf{Simulation / Engine} & Beat sequencing, pathology application, plugin stacking, time stepping &
  \file{graph\_engine.py}, \file{cable\_engine.py}, \file{simulation\_worker.py} \\
\hline
\textbf{Physics / Core} & Cellular models, conduction graph, cable propagation, ECG forward model &
  \file{fitzhugh\_nagumo.py}, \file{aliev\_panfilov.py}, \file{cable\_1d.py}, \file{lead\_field.py} \\
\hline
\end{tabularx}
\caption{Three-layer system architecture}
\label{tab:architecture_layers}
\end{table}

\subsection{1.3.2 Simulation Workflow}

The core simulation loop (executed at 500 Hz by the \class{SimulationWorker} thread) follows this sequence:

\begin{enumerate}
  \item \textbf{Parameter Update}: Check for parameter changes from GUI; apply base parameters and active pathologies to create \class{SimulationParameters}.
  \item \textbf{Engine Step}: Call \code{engine.step(dt)} with time increment $dt = 0.002$ s (500 Hz sample rate).
  \item \textbf{ECG Computation}: Engine computes 12 lead voltages from current activation state and returns \class{ECGSample}.
  \item \textbf{Display Update}: GUI thread updates scrolling ECG display with new sample.
  \item \textbf{Loop}: Repeat until user stops simulation.
\end{enumerate}

\note{The engine selection (Phase 1/2 vs Phase 3) is made via toolbar combo box; both engines implement the same \class{SimulationEngine} interface and are swappable at runtime.}

\section{1.4 Data Flow Diagram}

\begin{figure}[H]
\centering
\begin{tikzpicture}[
  node distance=2cm,
  block/.style={rectangle, draw, minimum width=2cm, minimum height=1cm, text centered},
  arrow/.style={->, thick},
]

% Boxes
\node[block] (gui) {GUI / Qt};
\node[block, right of=gui] (params) {Parameter Panel};
\node[block, right of=params] (pathology) {Pathology Panel};

\node[block, below of=gui] (worker) {Simulation Worker};
\node[block, right of=worker] (engine) {Simulation Engine\\(Graph or Cable)};
\node[block, right of=engine] (ecg) {ECG Display};

\node[block, below of=worker] (graph) {Conduction Graph\\(Phase 1/2)};
\node[block, right of=graph] (cable) {Cable 1D\\(Phase 3)};
\node[block, right of=cable] (fm) {Forward Model};

\node[block, below of=graph] (cellmodel) {Cell Models\\(FHN, AP)};
\node[block, right of=cellmodel] (leads) {Lead Vectors};

% Arrows
\draw[arrow] (gui) -- (params);
\draw[arrow] (params) -- (pathology);
\draw[arrow] (params.south) -- (worker.north);
\draw[arrow] (pathology.south) -- (worker.north);
\draw[arrow] (worker) -- (engine);
\draw[arrow] (engine) -- (ecg);

\draw[arrow] (engine.south) -- (graph.north);
\draw[arrow] (engine.south) -- (cable.north);
\draw[arrow] (graph) -- (fm);
\draw[arrow] (cable) -- (fm);

\draw[arrow] (graph.south) -- (cellmodel.north);
\draw[arrow] (cable.south) -- (cellmodel.north);
\draw[arrow] (fm.south) -- (leads.north);

\end{tikzpicture}
\caption{Data flow from user input through simulation to display}
\label{fig:dataflow}
\end{figure}

\section{1.5 Supported Pathologies (Phase 1/2)}

The Phase 1/2 engine natively supports the following pathological patterns via plugin architecture:

\begin{table}[H]
\centering
\begin{tabularx}{\textwidth}{|l|X|l|}
\hline
\textbf{Pathology} & \textbf{Mechanism} & \textbf{Type} \\
\hline
Sinus Bradycardia & Increased SA cycle length & Static \\
Sinus Tachycardia & Decreased SA cycle length & Static \\
AV Block I° & Prolonged AV node delay (PR \textgreater 200 ms) & Static \\
AV Block II° Mobitz I & Progressive PR lengthening (Wenckebach) & Stateful \\
AV Block II° Mobitz II & Fixed dropped beats & Stateful \\
AV Block III° & Complete AV dissociation (no QRS) & Static \\
LBBB & Left bundle branch conductance = 0 & Static \\
RBBB & Right bundle branch conductance = 0 & Static \\
LAHB & Left anterior fascicle block & Static \\
LPHB & Left posterior fascicle block & Static \\
Atrial Fibrillation & Suppress SA, rapid irregular QRS with f-waves & Stateful \\
Premature Atrial Contractions & Ectopic atrial firing & Stateful \\
Premature Ventricular Contractions & Ectopic ventricular firing & Stateful \\
\hline
\end{tabularx}
\caption{Supported pathological patterns in Phase 1/2 engine}
\label{tab:pathologies}
\end{table}

\note{Phase 3 engine currently supports AV blocks and sinus rate changes; bundle branch blocks fall back to Phase 1/2 morphology for visibility.}

% ============================================================================
% CHAPTER 2: CARDIAC ELECTROPHYSIOLOGY FROM FIRST PRINCIPLES
% ============================================================================

\chapter{Cardiac Electrophysiology — From First Principles}

This chapter builds cardiac electrophysiology systematically from biological mechanisms, through physical principles, to mathematical formulation. Each section assumes no prior knowledge and establishes concepts needed for later sections.

\section{2.1 Biological Foundations}

\subsection{2.1.1 The Cardiac Conduction System}

The heart beats because of a coordinated sequence of electrical activations that propagate through specialized tissues:

\begin{enumerate}
  \item \textbf{Sinoatrial (SA) Node} ($\approx$ 60 -- 100 action potentials per minute):
    \begin{itemize}
      \item Located at the junction of the superior vena cava and right atrium.
      \item Contains \emph{pacemaker cells} that spontaneously depolarize (funny current $I_f$ via HCN channels).
      \item Fires once per heartbeat, initiating the next cardiac cycle.
      \item Primary rate-determining site.
    \end{itemize}

  \item \textbf{Atrial Myocardium} (atrial muscle):
    \begin{itemize}
      \item Conduction velocity $\approx 0.5 -- 1.0$ m/s (slower than ventricular).
      \item Activation spreads right-to-left and inferiorly.
      \item Produces the P wave on ECG (depolarization only; repolarization hidden under QRS).
      \item Refractory period $\approx 150 -- 200$ ms.
    \end{itemize}

  \item \textbf{Atrioventricular (AV) Node}:
    \begin{itemize}
      \item Located between atria and ventricles, at the junction of the atrial septum.
      \item Narrow conduction velocity ($\approx 0.05$ m/s) and long conduction delay ($\approx$ 100 ms from atrial entry to His bundle exit).
      \item \emph{Physiological gating}: allows time for atrial contraction to complete before ventricular activation.
      \item Refractory period $\approx$ 300 ms; effectively blocks beats arriving too early.
    \end{itemize}

  \item \textbf{His Bundle \& Bundle Branches}:
    \begin{itemize}
      \item Fast-conducting Purkinje tissue ($\approx 2 -- 4$ m/s).
      \item Divides into left bundle branch (LBB) and right bundle branch (RBB).
      \item Conducts quickly to both ventricles so they depolarize nearly synchronously.
    \end{itemize}

  \item \textbf{Ventricular Myocardium}:
    \begin{itemize}
      \item Working myocardium with conduction velocity $\approx 0.4 -- 0.5$ m/s.
      \item Activated by Purkinje terminal branches.
      \item Left ventricle usually activates before right (leftward bias).
      \item Base-to-apex activation gradient due to apical Purkinje penetration.
      \item Refractory period $\approx 250 -- 300$ ms.
    \end{itemize}
\end{enumerate}

\subsection{2.1.2 The Cardiac Action Potential}

Every cardiac myocyte (muscle cell) and pacemaker cell undergoes a cyclical electrical event called the \concept{action potential} (AP). The AP is divided into five phases:

\begin{figure}[H]
\centering
\begin{tikzpicture}[scale=0.8]
  \begin{axis}[
    xlabel={Time (ms)},
    ylabel={Membrane Potential (mV)},
    xmin=0, xmax=400,
    ymin=-100, ymax=50,
    grid=both,
    grid style={draw=gray!30},
  ]
    % Phase 0: rapid depolarization
    \addplot[very thick, blue, no markers] coordinates {(0, -85) (5, 30)};
    % Phase 1: early repolarization
    \addplot[very thick, blue, no markers] coordinates {(5, 30) (15, 5)};
    % Phase 2: plateau
    \addplot[very thick, blue, no markers] coordinates {(15, 5) (250, 0)};
    % Phase 3: rapid repolarization
    \addplot[very thick, blue, no markers] coordinates {(250, 0) (280, -85)};
    % Phase 4: resting
    \addplot[very thick, blue, no markers] coordinates {(280, -85) (400, -85)};
    
    % Labels
    \node at (2.5, 0) {0};
    \node at (10, 20) {1};
    \node at (130, 5) {2};
    \node at (265, -40) {3};
    \node at (350, -87) {4};
  \end{axis}
\end{tikzpicture}
\caption{Five phases of the ventricular cardiac action potential}
\label{fig:ap_phases}
\end{figure}

The five phases are:

\begin{enumerate}
  \item \textbf{Phase 0 (Depolarization, 0--5 ms)}: Rapid inward sodium current ($I_{Na}$) driven by Na$^+$ gradient (outside positive, inside negative). Membrane potential rises from $-85$ mV to $+30$ mV in $\approx$ 5 ms.
    \begin{itemize}
      \item Sodium channels open (voltage-gated, opening at $-60$ mV).
      \item Potassium channels close (inactivate).
      \item Creates the characteristic steep upstroke on ECG.
    \end{itemize}

  \item \textbf{Phase 1 (Early Repolarization, 5--15 ms)}: Brief outward current from potassium (transient outward $I_{to}$), producing a slight "notch" on the upstroke. Current decreases as inactivation dominates.

  \item \textbf{Phase 2 (Plateau, 15--250 ms)}: Balance between inward L-type calcium current ($I_L$) and outward potassium current ($I_K$). Membrane potential held near 0 mV for a prolonged period.
    \begin{itemize}
      \item Calcium enters cell, triggering contraction.
      \item Duration determines the action potential duration (APD), typically 250--300 ms in ventricular myocytes.
    \end{itemize}

  \item \textbf{Phase 3 (Repolarization, 250--280 ms)}: Delayed rectifier potassium channels ($I_K$) dominate; inward currents decline. Membrane potential returns toward $-85$ mV.
    \begin{itemize}
      \item Generates the T wave on ECG.
      \item Duration determines when the cell is excitable again (after full repolarization).
    \end{itemize}

  \item \textbf{Phase 4 (Resting, 280+ ms)}: Cell returns to resting potential ($-85$ mV). Membrane is permeable to potassium (maintaining resting potential via Na$^+$/K$^+$-ATPase).
    \begin{itemize}
      \item In pacemaker cells (SA/AV node), spontaneous diastolic depolarization occurs via $I_f$.
      \item In myocytes, potential is stable until the next stimulus arrives.
    \end{itemize}
\end{enumerate}

\subsection{2.1.3 Refractoriness}

The refractory period defines when a cell can fire again after an AP:

\begin{itemize}
  \item \textbf{Absolute Refractory Period (ARP, $\approx 250$ ms)}: Sodium channels are inactivated (cannot open regardless of stimulus). No new AP can be generated, no matter how strong the stimulus.
  
  \item \textbf{Relative Refractory Period (RRP, $\approx 50 -- 100$ ms after ARP)}: Sodium channels have recovered from inactivation, but membrane potential is still negative relative to resting. A \emph{supranormal} stimulus (stronger than usual) can evoke an AP, but conduction will be slow and amplitude reduced.
\end{itemize}

In the simulator, refractoriness is:
- \textbf{Phase 1/2}: Hardcoded per-node values (e.g., ventricular $\approx$ 250 ms).
- \textbf{Phase 3}: \emph{Emergent} from FHN dynamics; refractory state arises naturally from slow recovery dynamics.

\subsection{2.1.4 Conduction Velocity}

AP propagation speed varies by tissue type:

\begin{table}[H]
\centering
\begin{tabularx}{\textwidth}{|l|c|X|}
\hline
\textbf{Tissue} & \textbf{Velocity (m/s)} & \textbf{Physiological Basis} \\
\hline
SA node & $\approx 0.05$ & Calcium-dependent, sparse gap junctions, slow channel kinetics \\
Atrial myocardium & 0.5 -- 1.0 & Moderate Na$^+$ channel density, good gap junctions \\
AV node & $\approx 0.05$ & Calcium-dependent, narrow cross-section, slow conduction \\
His bundle / Purkinje & 2 -- 4 & Very high Na$^+$ channel density, abundant gap junctions, large diameter \\
Ventricular myocardium & 0.4 -- 0.5 & Working myocardium, good but not optimal channel density \\
\hline
\end{tabularx}
\caption{Conduction velocities in different cardiac tissues}
\label{tab:conduction_velocities}
\end{table}

\subsubsection{Physical Basis of Conduction Velocity}

Conduction velocity is determined by the cable equation (see \S\ref{sec:cable_equation}):

$$v_c = \sqrt{\frac{g \cdot D}{2 \cdot C_m \cdot r_i}}$$

where:
- $g$ = conductance per unit length (related to $I_{Na}$ availability and gap junction conductance)
- $D$ = diffusion coefficient (related to gap junction resistance)
- $C_m$ = membrane capacitance per unit area
- $r_i$ = intracellular resistance per unit length

Tissues with large diameter fibers and abundant gap junctions (like Purkinje) have low $r_i$, yielding fast conduction. Tissues with poor gap junctions (like AV node) have high $r_i$, yielding slow conduction.

\section{2.2 Electrical Principles: The ECG Forward Problem}

\subsection{2.2.1 Dipole Approximation}

The heart can be modeled as a moving electrical dipole vector $\mathbf{p}(t)$ in 3D space. During depolarization and repolarization, different regions of the heart become active; the net dipole is the vector sum of all regional contributions.

The electrical potential at a distant point due to a dipole is given by:

$$\Phi(\mathbf{r}) = \frac{1}{4\pi\sigma} \frac{\mathbf{p} \cdot \mathbf{r}}{r^3}$$

where:
- $\mathbf{p}$ = dipole moment (units: \textmu A·cm or equivalent)
- $\mathbf{r}$ = position vector from dipole to observation point
- $r = |\mathbf{r}|$ = distance
- $\sigma$ = torso conductivity

The ECG records the potential difference between two electrodes. For a single lead with lead vector $\mathbf{L}$ (unit vector from negative to positive electrode):

$$V_{\text{lead}} = \mathbf{L} \cdot \nabla\Phi = \frac{1}{4\pi\sigma} \frac{\mathbf{L} \cdot \mathbf{p}}{r^2}$$

If we normalize by assuming $\frac{1}{4\pi\sigma r^2}$ is constant, the lead voltage is simply the projection of the dipole onto the lead vector:

$$V_{\text{lead}} \propto \mathbf{L} \cdot \mathbf{p}$$

\subsection{2.2.2 The 12-Lead ECG}

Standard clinical ECG records voltage from 12 locations (leads):

\begin{enumerate}
  \item \textbf{Frontal plane leads (I, II, III, aVR, aVL, aVF)}:
    \begin{itemize}
      \item Lie in the frontal plane perpendicular to the body's long axis.
      \item Record the component of the dipole in the frontal plane.
      \item Einthoven's triangle: Lead I (0°), II (+60°), III (+120°), aVR (−150°), aVL (−30°), aVF (+90°).
    \end{itemize}

  \item \textbf{Precordial leads (V1--V6)}:
    \begin{itemize}
      \item Lie on the chest around the heart, roughly in the transverse plane.
      \item V1 (right anterior) to V6 (left lateral): "sees" the progression of activation from right to left.
      \item Useful for detecting localized abnormalities (e.g., anterior MI, lateral infarction).
    \end{itemize}
\end{enumerate}

\subsection{2.2.3 Lead Vectors in the Simulator}

The simulator stores lead vectors as unit vectors $\mathbf{L}_i$ (one per lead) in 3D Cartesian coordinates:

$$\mathbf{L}_i = \begin{pmatrix} L_{i,x} \\ L_{i,y} \\ L_{i,z} \end{pmatrix}$$

where:
- $x$: leftward (patient's left)
- $y$: inferior (feet-ward)
- $z$: anterior (front of chest)

The 12-lead voltage vector is:

$$\mathbf{V}_{\text{ECG}} = \mathbf{L} \cdot \mathbf{p}_{\text{total}}(t)$$

where $\mathbf{L}$ is the $(12 \times 3)$ lead matrix (rows are lead vectors) and $\mathbf{p}_{\text{total}}(t)$ is the instantaneous total cardiac dipole.

\section{2.3 Mathematical Models of Excitability}

\subsection{2.3.1 FitzHugh-Nagumo Model}

The FitzHugh-Nagumo (FHN) model is a 2-variable simplified model that captures the essential qualitative behavior of cardiac action potentials without the complexity of full ion-channel models.

\subsubsection{Equations (Dimensionless Form)}

In dimensionless time $\tau$:

\begin{equation}
\frac{dv}{d\tau} = v - \frac{v^3}{3} - w + I_{\text{ext}}
\label{eq:fhn_v}
\end{equation}

\begin{equation}
\frac{dw}{d\tau} = \varepsilon (v + a - b \cdot w)
\label{eq:fhn_w}
\end{equation}

where:
- $v$ = fast variable (membrane potential analogue); dimensionless, ranges $\approx [-1.2, 2.0]$
- $w$ = slow variable (recovery variable); dimensionless, ranges $\approx [-0.625, 0.2]$
- $\varepsilon$ = time-scale separation parameter (small, $\approx 0.08$)
- $a, b$ = model parameters (typically $a \approx 0.7$, $b \approx 0.8$)
- $I_{\text{ext}}$ = external stimulus current (dimensionless)

\subsubsection{Physical Time Conversion}

The dimensionless equations are scaled to physical time via a time constant $\tau_s$ (seconds):

\begin{equation}
\frac{dv}{dt} = \frac{1}{\tau_s}\left(v - \frac{v^3}{3} - w + I_{\text{ext}}\right)
\end{equation}

\begin{equation}
\frac{dw}{dt} = \frac{\varepsilon}{\tau_s}(v + a - b \cdot w)
\end{equation}

In the simulator, $\tau_s = 0.015$ s = 15 ms (calibrated to match clinical PR and QRS intervals; see \S\ref{sec:phase3_calibration}).

\subsubsection{Phase Plane Analysis}

The two nullclines (curves where $dv/d\tau = 0$ and $dw/d\tau = 0$) determine the qualitative behavior:

\begin{itemize}
  \item \textbf{$v$-nullcline}: $w = v - \frac{v^3}{3} + I_{\text{ext}}$ (cubic, N-shaped when $I=0$).
  \item \textbf{$w$-nullcline}: $w = \frac{v + a}{b}$ (linear).
\end{itemize}

The fixed point (resting state) is the intersection of the nullclines. For $a = 0.7, b = 0.8, \varepsilon = 0.08, I_{\text{ext}} = 0$:

$$v_{\text{rest}} \approx -1.200, \quad w_{\text{rest}} \approx -0.625$$

When stimulated with $I_{\text{ext}} > 0$:

\begin{enumerate}
  \item $v$ rapidly increases (fast variable; $\varepsilon$ is small).
  \item $w$ slowly increases (recovery; slow variable).
  \item When $I_{\text{ext}}$ is turned off, $v$ repolarizes rapidly (returns to rest).
  \item $w$ repolarizes slowly (takes $\approx 200$ ms), creating the refractory period.
\end{enumerate}

\note{The FHN model is \emph{not} quantitatively accurate for specific ionic currents but captures the essential dynamics: excitability, AP shape, and emergent refractoriness.}

\subsubsection{Action Potential Duration (APD)}

In the FHN model, the APD is determined primarily by the recovery time scale:

$$\text{APD}_{90} \approx 2.3 \times \frac{\tau_s}{b \cdot \varepsilon}$$

With $\tau_s = 0.015$ s, $b = 0.8$, $\varepsilon = 0.08$:

$$\text{APD}_{90} \approx 2.3 \times \frac{0.015}{0.8 \times 0.08} \approx 0.54 \text{ s} = 540 \text{ ms}$$

This is longer than clinical (250--300 ms) and will be refined in later phases.

\subsubsection{Effective Refractory Period (ERP)}

The ERP in FHN emerges naturally: when $w$ is elevated (slow recovery), $v$ cannot cross the threshold to fire again. The ERP duration is approximately:

$$\text{ERP} \approx \frac{2}{\varepsilon \cdot b \cdot \tau_s} \approx \frac{2}{0.08 \times 0.8 \times 0.015} \approx 208 \text{ ms}$$

At a 70 bpm heart rate (RR = 857 ms), the cell repolarizes and recovers fully before the next beat, allowing normal rhythmic firing.

\subsection{2.3.2 Aliev-Panfilov Model}

The Aliev-Panfilov (AP) model is another 2-variable model designed specifically for cardiac tissue, with improved plateau dynamics:

\begin{equation}
\frac{du}{d\tau} = -k \cdot u(u - a)(u - 1) - u \cdot v + I_{\text{ext}}
\end{equation}

\begin{equation}
\frac{dv}{d\tau} = \left(\varepsilon_0 + \frac{\mu_1 \cdot v}{\mu_2 + u}\right) \cdot (-v - k \cdot u(u - a - 1))
\end{equation}

where:
- $u$ = membrane potential ($0$ = rest, $1$ = depolarized)
- $v$ = recovery variable
- $k, a$ = excitability parameters
- $\varepsilon_0, \mu_1, \mu_2$ = recovery rate and adaptivity coefficients

The key feature is the adaptive recovery time scale $\varepsilon(u, v)$, which accelerates repolarization and creates a physiologically realistic plateau. However, the FHN model is simpler and currently used in Phase 3; Aliev-Panfilov is available for future use.

\subsection{2.3.3 Stimulus-Response Curves}

For both FHN and AP models, there exists a minimum stimulus strength (threshold) required to evoke an AP:

\begin{equation}
I_{\text{threshold}} = \frac{2v_{\text{rest}}^3 + 3v_{\text{rest}}^2}{3(1 + \text{time scale factor})}
\end{equation}

For FHN: $I_{\text{threshold}} \approx 0.4$ dimensionless units.

In practice, the simulator uses stimulus amplitude of $1.5$ (well above threshold) to ensure reliable AP initiation in SA pacemaker.

% ============================================================================
% CHAPTER 3: PHASE 1/2 CONDUCTION GRAPH ENGINE
% ============================================================================

\chapter{Phase 1/2: Discrete Conduction Graph Engine}

This chapter describes the Phase 1/2 simulation engine, which models cardiac activation as a discrete directed acyclic graph (DAG) with hardcoded conduction times and refractory periods.

\section{3.1 Conduction Graph Architecture}

\subsection{3.1.1 Node Definition}

Each cardiac region is represented as a single \emph{node} in an abstract graph. The simulator defines 18 nodes:

\begin{table}[H]
\centering
\begin{tabularx}{\textwidth}{|l|X|X|}
\hline
\textbf{Node Name} & \textbf{Anatomical Region} & \textbf{Default Activation Time} \\
\hline
\multicolumn{3}{|c|}{\textbf{Atrial Nodes}} \\
\hline
SA\_NODE & Sinoatrial node (pacemaker) & 0 ms (reference) \\
RA & Right atrium & $\approx$ 45 ms \\
LA & Left atrium & $\approx$ 90 ms \\
\hline
\multicolumn{3}{|c|}{\textbf{AV Conduction}} \\
\hline
AV\_NODE & Atrioventricular node & $\approx$ 115 ms \\
\hline
\multicolumn{3}{|c|}{\textbf{Ventricular}} \\
\hline
HIS & His bundle & $\approx$ 135 ms \\
LBB & Left bundle branch & $\approx$ 145 ms \\
RBB & Right bundle branch & $\approx$ 145 ms \\
SEPT\_EARLY & Early septal activation & $\approx$ 155 ms \\
SEPT\_MID & Mid-septal activation & $\approx$ 165 ms \\
RV & Right ventricle & $\approx$ 175 ms \\
LV\_ANT & LV anterior wall & $\approx$ 175 ms \\
LV\_LAT & LV lateral wall & $\approx$ 200 ms \\
LV\_INF & LV inferior wall & $\approx$ 205 ms \\
LV\_POST & LV posterior wall & $\approx$ 215 ms \\
\hline
\multicolumn{3}{|c|}{\textbf{Repolarization}} \\
\hline
(Implicit from node activation + APD) & T wave on ECG & QRS end $+$ $\approx$ 280 ms \\
\hline
\end{tabularx}
\caption{ECG nodes in Phase 1/2 graph engine}
\label{tab:ecg_nodes}
\end{table}

\subsection{3.1.2 Graph Edges and Conduction Delays}

The graph is defined as a set of directed edges, each with a conduction delay:

\begin{itemize}
  \item \textbf{Forward (normal) edges}: SA → RA, RA → LA, LA → AV, AV → HIS, HIS → LBB, HIS → RBB, LBB → SEPT\_EARLY, etc.
  \item \textbf{Retrograde edges}: Added in Phase 2 to model bundle branch blocks. For example, if LBB is blocked, RV can activate LV retrogradely through working myocardium.
\end{itemize}

When a node is activated at time $t_i$, all its downstream nodes are scheduled for activation at times:

$$t_j = t_i + \text{delay}_{i \to j}$$

subject to refractory period constraints (see \S\ref{sec:refractoriness}).

\subsection{3.1.3 \class{ConductionGraph} Class}

\textbf{Location}: \file{cardiac\_sim/core/conduction/graph.py}

\textbf{Responsibility}: Represents the cardiac conduction DAG and computes activation times via breadth-first search (BFS).

\textbf{Key Method}: \code{run\_bfs(sa\_fire\_time, params) → dict[str, float]}

Computes absolute activation times for all nodes after SA fires at \code{sa\_fire\_time}. Returns dict mapping node name → activation time (seconds).

Algorithm:

\begin{enumerate}
  \item Initialize queue with SA\_NODE at time 0.
  \item For each node in queue:
    \begin{itemize}
      \item Get all outgoing edges (predecessors in DAG).
      \item For each successor:
        \begin{itemize}
          \item Compute activation time: $t_j = t_i +$ delay.
          \item Check refractory period: if $t_j < t_{\text{last\_activation}}[j] +$ ERP, skip (refractory).
          \item Otherwise, queue for activation and record time.
        \end{itemize}
    \end{itemize}
  \item Return final activation times dict.
\end{enumerate}

\textbf{Complexity}: $O(V + E)$ where $V$ = number of nodes ($\approx 18$) and $E$ = number of edges ($\approx 30$).

\section{3.2 Refractoriness in Phase 1/2}

\subsection{3.2.1 Hardcoded Refractory Periods}

In Phase 1/2, each node has a fixed effective refractory period (ERP):

\begin{code}
cardiac_sim/core/conduction/graph.py (lines 50--80):

REFRACTORY_PERIODS = {
    "SA_NODE":     100.0e-3,   # 100 ms
    "RA":          200.0e-3,   # 200 ms
    "LA":          200.0e-3,   # 200 ms
    "AV_NODE":     300.0e-3,   # 300 ms
    "HIS":         250.0e-3,   # 250 ms
    "LBB":         250.0e-3,   # 250 ms
    "RBB":         250.0e-3,   # 250 ms
    "SEPT_EARLY":  250.0e-3,   # 250 ms
    "SEPT_MID":    250.0e-3,   # 250 ms
    "RV":          250.0e-3,   # 250 ms
    "LV_ANT":      250.0e-3,   # 250 ms
    "LV_LAT":      250.0e-3,   # 250 ms
    "LV_INF":      250.0e-3,   # 250 ms
    "LV_POST":     250.0e-3,   # 250 ms
}
\end{code}

These values are fixed and represent a typical physiological baseline. They are overridden by \class{SimulationParameters} for plugin use (e.g., AV block plugins modify AV node delay and conductance).

\subsection{3.2.2 Refractory Logic}

During BFS, when a candidate activation arrives at node $j$ at time $t_j$:

\begin{enumerate}
  \item \textbf{Check if node has been activated before}: If not in \code{last\_activation} dict, activate immediately.
  \item \textbf{Check refractory period}: If $t_j < \text{last\_activation}[j] +$ ERP$[j]$, the node is refractory → skip activation.
  \item \textbf{Otherwise}: Activate and update \code{last\_activation}[j] = $t_j$.
\end{enumerate}

\textbf{Effect}: A rapidly-firing stimulus (e.g., premature contraction) will fail to propagate past a refractory node, stopping the wave and preventing retrograde conduction.

\section{3.3 ECG Computation in Phase 1/2}

\subsection{3.3.1 Dipole Contribution Per Node}

Each node is associated with a dipole contribution $\mathbf{p}_i(t)$ that evolves in time:

\begin{equation}
\mathbf{p}_i(t) = A_{\text{depol}} \cdot \exp\left(-\frac{(t - t_{\text{act},i})^2}{2\sigma_{\text{depol}}^2}\right) \mathbf{d}_i
+ A_{\text{repol}} \cdot \exp\left(-\frac{(t - t_{\text{repol},i})^2}{2\sigma_{\text{repol}}^2}\right) \mathbf{d}_i
\label{eq:dipole_gaussian}
\end{equation}

where:
- $t_{\text{act},i}$ = absolute activation time of node $i$
- $t_{\text{repol},i} = t_{\text{act},i} +$ APD (repolarization time)
- $\sigma_{\text{depol}}, \sigma_{\text{repol}}$ = Gaussian widths (represent the duration over which the dipole is "smeared")
- $A_{\text{depol}}, A_{\text{repol}}$ = depolarization and repolarization amplitudes (typically $A_{\text{repol}} \approx -0.3 \times A_{\text{depol}}$)
- $\mathbf{d}_i$ = dipole direction vector (fixed per node)

\subsection{3.3.2 Total Dipole and Lead Voltages}

The total instantaneous dipole at time $t$ is the sum over all nodes in currently active beats:

\begin{equation}
\mathbf{p}_{\text{total}}(t) = \sum_{i \in \text{active\_nodes}} \mathbf{p}_i(t)
\end{equation}

The 12-lead voltages are:

\begin{equation}
V_k(t) = \mathbf{L}_k \cdot \mathbf{p}_{\text{total}}(t), \quad k = 1, \ldots, 12
\end{equation}

where $\mathbf{L}_k$ is the $k$-th lead vector.

\textbf{Implementation}: \code{cardiac\_sim/simulation/graph\_engine.py}, method \code{\_compute\_ecg()}.

\section{3.4 Plugin Architecture (Static vs Stateful)}

\subsection{3.4.1 Static Plugins}

Static plugins are applied once when the user selects them, and their effect is incorporated into the cached parameter set for all subsequent beats.

\textbf{Example: LBBB plugin}

\begin{lstlisting}
class LBBB(AbstractPathologyPlugin):
    def apply(self, params: SimulationParameters) -> SimulationParameters:
        p = params.copy()
        p.his_bundle.left_branch_conductance = 0.0  # Block LBB
        return p
\end{lstlisting}

When the graph is rebuilt with this parameter set, the forward edge HIS → LBB is disabled (conductance = 0 means no conduction). Activation must find an alternate path; the retrograde working-myocardium edge activates LV before RV, producing the characteristic wide QRS morphology.

\textbf{Lifecycle}:

\begin{enumerate}
  \item User selects LBBB from pathology panel.
  \item GUI calls \code{engine.apply\_pathology(LBBB())}
  \item Engine locks, calls \code{plugin.apply(params)} → modified params.
  \item Engine rebuilds graph with new params.
  \item All subsequent beats use the modified graph.
  \item Graph is not rebuilt until another static plugin is applied or removed.
\end{enumerate}

\subsection{3.4.2 Stateful Plugins}

Stateful plugins are called \emph{once per beat} and can make beat-specific parameter overrides that don't persist.

\textbf{Example: Wenckebach (AV Block II° Mobitz I) plugin}

\begin{lstlisting}
class AVBlockMobitzI(AbstractStatefulPathologyPlugin):
    def __init__(self):
        self._position = 0  # Position in Wenckebach cycle
        self._cycle_length = 4  # 4 P waves → 3 QRS (4:3 conduction)
    
    def on_beat(self, params, beat_id):
        p = params.copy()
        pos = self._position % self._cycle_length
        if pos == self._cycle_length - 1:
            # Dropped beat: complete block
            p.av_node.conductance = 0.0
        else:
            # Conducted beat: PR lengthens progressively
            p.av_node.conduction_delay_ms = 160 + pos * 60
            p.av_node.conductance = 1.0
        self._position += 1
        return p
\end{lstlisting}

\textbf{Lifecycle}:

\begin{enumerate}
  \item User selects Wenckebach from pathology panel.
  \item GUI calls \code{engine.apply\_pathology(AVBlockMobitzI())} (registered as stateful).
  \item Each beat:
    \begin{enumerate}
      \item Engine calls \code{plugin.on\_beat(cached\_params, beat\_id)}.
      \item Plugin returns beat-specific params (e.g., conductance = 0 for dropped beat).
      \item Temporary graph is built from beat-specific params.
      \item BFS run computes activation times for this beat only.
      \item After beat completes, temporary graph is discarded.
      \item Internal state (e.g., \code{\_position}) is preserved for next beat.
    \end{enumerate}
  \item When plugin is removed, internal state is reset via \code{reset()}.
\end{enumerate}

% [Document continues with Chapters 4-8...]
% [Due to length, I'm showing the structure and will continue in the next message]

\end{document}
